\section{The slice \texorpdfstring{$j_1=\frac32$}{j1 = 3/2}}\label{app:z32}

This appendix proves \cref{thm:z32}. The proof has four parts: the reduction, the chain
lemmas with their $30$ polynomial certificates, a finite box of $6.4\times10^7$ parameter
sets evaluated in outward-rounded interval arithmetic (\cref{subsec:box32}), and the tail
lemmas with their $50$ certificates.

\subsection{Zigzags and links}

Fix an ordered pairing with pairs $(n_1,m_1)$, $(n_2,m_2)$ and write
$P_i(L)=(n_i^2-L^2)(L^2-m_i^2)$. For $s=\frac32$, $F(x)=4-x^2$ and
$G(u)=(u+\frac52)(u-\frac32)$. Zigzag $0$ has rows $R_b=(b,b+\frac32)$ and columns
$C_b=(b,b+\frac12)$; zigzag $1$ has rows $R'_b=(b,b-\frac12)$ and columns $C'_b=(b,b-\frac32)$.
In the order of $u$ each zigzag is a path $C_b,R_b,C_{b+1},\dots$, and the block is
$B=\bigl(\begin{smallmatrix}D_0&0\\N&D_1\end{smallmatrix}\bigr)$. By \cref{prop:grid} the squared
links of $T$ are
\begin{center}\small
\begin{tabular}{lll}
\toprule
\estabhead{link} & \estabhead{between} & \estabhead{squared weight}\\
\midrule
$q_b$ & $C_b$--$R_b$ & $\frac3{16}P_2(b+\frac32)\,b/(b+1)$\\
$p_b$ & $R_b$--$C_{b+1}$ & $\frac3{16}P_1(b+1)\,(2b+5)/(2b+3)$\\
$q'_b$ & $C'_b$--$R'_b$ & $\frac3{16}P_2(b-\frac12)\,(b+1)/b$\\
$p'_b$ & $R'_b$--$C'_{b+1}$ & $\frac3{16}P_1(b+1)\,(2b-1)/(2b+1)$\\
$n^+_b$ & $R'_b$--$C_b$ & $P_2(b+\frac12)(2b+3)(2b-1)/(16b(b+1))$\\
$n^-_b$ & $R'_b$--$C_{b-1}$ & $P_1(b)(b^2-1)/(4b^2-1)$\\
\bottomrule
\end{tabular}
\end{center}
(checked exactly against the grid on $4\,171$ links). The site ranges follow from the
region: rows of zigzag $0$ have $b\in[\max(m_1,m_2-\frac32,0),\min(n_1-1,n_2-\frac52)]$,
its columns $b\in[\max(m_1,m_2-\frac12,\frac12),\min(n_1-1,n_2-\frac32)]$, rows of zigzag $1$
$b\in[\max(m_1,m_2+\frac12,1),\min(n_1-1,n_2-\frac12)]$, and its columns
$b\in[\max(m_1,m_2+\frac32,\frac32),\min(n_1-1,n_2+\frac12)]$. A zigzag is of type CR, RC,
RR or CC according to whether its path starts with a column or a row and ends with a
column or a row; its imbalance is $0$, $0$, $+1$ or $-1$.

\subsection{Reduction}

\begin{lemma}[All-tall pairings]\label{lem:alltall}
If in some ordered pairing both zigzags have imbalance $0$ or $+1$, then $B$ is injective
and $\lambda(T)=\smin(B)$.
\end{lemma}

\begin{proof}
Each $D_k$ is a path with nonzero links and at least as many rows as columns, so it is
injective (delete an end row if it is tall). If $Bx=0$ then $D_0x_0=0$, so $x_0=0$, and then
$D_1x_1=0$.
\end{proof}

The proof needs a pairing in which neither zigzag has more columns than rows, and which
pairing works depends on the gaps between the sorted spins of legs $2$--$5$. The
regimes below sort the tuples by those gaps. In R1 and R4 one outer gap exceeds the
other by at least $\frac32$ ($g_1$ in R1, $g_3$ in R4); in R2$^\mp$ they differ by
exactly $\frac12$; R1c and R4c are the boundary cases in which the smaller outer gap
vanishes; and Tm, Tp are the two one-parameter families in which only one leg differs
from the others, by $-\frac12$ (Tm, the smallest leg) or $+\frac12$ (Tp, the largest).

\begin{lemma}[Regimes]\label{lem:regimes}
Sort $j_2\le j_3\le j_4\le j_5$ and put $g_1=j_3-j_2$, $g_2=j_4-j_3$, $g_3=j_5-j_4$; then
$g_1+g_3\equiv\frac12\pmod1$. Every tuple falls in exactly one of the following regimes,
and in the pairing listed both zigzags have imbalance in $\{0,1\}$.
\begin{center}\small
\begin{tabular}{llll}
\toprule
\estabhead{regime} & \estabhead{gaps} & \estabhead{pairing} & \estabhead{zigzag 0 / zigzag 1}\\
\midrule
R1 & $g_1-g_3\ge\frac32$ & $(j_2j_4|j_3j_5)$ & CR / CR\\
R2$^-$ & $g_1-g_3=\frac12$, $g_3\ge\frac12$ & $(j_2j_4|j_3j_5)$ & CR / RR\\
R1c & $g_1=\frac12$, $g_3=0$, $g_2\ge\frac12$ & $(j_2j_3|j_4j_5)$ & CR / CR\\
Tm & $g_1=\frac12$, $g_2=g_3=0$ & $(j_3j_4|j_2j_5)$ & RC / RR\\
R4 & $g_3-g_1\ge\frac32$ & $(j_2j_4|j_3j_5)$ & RR / RR\\
R2$^+$ & $g_3-g_1=\frac12$, $g_1\ge\frac12$ & $(j_2j_4|j_3j_5)$ & CR / RR\\
R4c & $g_1=0$, $g_3=\frac12$, $g_2\ge\frac12$ & $(j_2j_3|j_4j_5)$ & RR / RR\\
Tp & $g_1=0$, $g_3=\frac12$, $g_2=0$ & $(j_2j_3|j_4j_5)$ & RC / RR\\
\bottomrule
\end{tabular}
\end{center}
The one-parameter families are Tm $=(\frac32;j,j+\frac12,j+\frac12,j+\frac12)$ and
Tp $=(\frac32;j,j,j,j+\frac12)$.
\end{lemma}

\begin{proof}
In $(j_2j_4|j_3j_5)$, $n_2-n_1=g_1+g_3$ and $m_2-m_1=g_3-g_1$, with $m_1=g_1+g_2$,
$m_2=g_2+g_3$; in $(j_2j_3|j_4j_5)$, $n_2-n_1=g_1+2g_2+g_3$ and $(m_1,m_2)=(g_1,g_3)$. The
conditions are disjoint and exhaustive because $g_1-g_3$ is a nonzero half-odd integer.
Substituting into the four range formulas gives the types; for example in R1,
$m_2+\frac32\le m_1$ and $m_1\ge\frac32$ make every lower end $m_1$, and $n_2\ge n_1+\frac32$
makes every upper end $n_1-1$.
\end{proof}

The census of all $54\,019$ tuples with spins at most $20$ (legs $2$--$5$ sorted,
$\dim\Inv\ge2$) confirms the types, the
imbalances and the ratio hypotheses below in exact arithmetic.

\begin{proposition}{The coupling identity}{coupling32}
On the columns of zigzag $0$, $N^TN=\frac43D_0^TD_0+\diag(\Lambda)$ with
$\Lambda_c=n^+_c{}^2+n^-_{c+1}{}^2-\frac43(q_c^2+p_{c-1}^2)=\Lambda^1_c+\Lambda^2_c$,
\[
  \Lambda^1_c=\frac{P_1(c+1)\,c(c+2)}{(2c+1)(2c+3)}-\frac{P_1(c)(2c+3)}{4(2c+1)},\qquad
  \Lambda^2_c=\frac{P_2(c+\frac12)(2c+3)(2c-1)}{16c(c+1)}-\frac{P_2(c+\frac32)\,c}{4(c+1)} .
\]
\end{proposition}

This is \cref{prop:zigzag} with $\mathsf r=F(0)/F(-1)=\frac43$; the formula holds at every
column, the ends included, because the link of a missing row vanishes.

\begin{corollary}\label{cor:FW}
Let $F_0\ge1/\lambda_{\min}(D_0^TD_0)$, $F_1\ge1/\lambda_{\min}(D_1^TD_1)$ and
$C_\Lambda \ge\sup_xx^T\Lambda x/|D_0x|^2$ ($C_\Lambda =0$ if $\Lambda\le0$). Then for every $\tau>0$,
\[
  \smin(B)^2\ge\min\Bigl(\frac{1-\tau(\frac43+C_\Lambda)}{F_0},\ \frac{\tau}{(1+\tau)F_1}\Bigr).
\]
When $D_0$ is square one may take $C_\Lambda =\sum_c\Lambda_c^+R_c$ with $R_c\ge\sum_k(D_0^{-1})_{ck}^2$.
\end{corollary}

\begin{proof}
If $1-\tau(\frac43+C_\Lambda)\le0$, the first term of the minimum is at most $0$ and there is
nothing to prove. Otherwise, by \eqref{eq:twochain} and $D_0^TD_0-\tau N^TN=(1-\frac43\tau)D_0^TD_0-\tau\Lambda\succeq
(1-\tau(\frac43+C_\Lambda))D_0^TD_0$; for the last statement write $y=D_0x$ and bound the operator
norm of $\Lambda^{+1/2}D_0^{-1}$ by its Frobenius norm.
\end{proof}

\subsection{Chains}

Each lemma uses only the explicit inverse of a bidiagonal path. For a row $R_b$ write
$r_b=p_b^2/q_b^2$ (same-row ratio) and $\bar r_b=p_b^2/q_{b+1}^2$ (same-column ratio).

\begin{lemma}[CR chains]\label{lem:CR}
If $D$ is square with $(Dx)_{R_b}=q_bx_b+p_bx_{b+1}$, then for $k\ge c$
\[
  |(D^{-1})_{ck}|=q_c^{-1}\prod_{t=c}^{k-1}\sqrt{\bar r_t}.
\] If every $\bar r_t\le1$, then
$1/\lambda_{\min}(D^TD)\le\sum_c(\mathrm{hi}-c+1)/q_c^2$ and $R_c\le(\mathrm{hi}-c+1)/q_c^2$
(count form). With $\ell_c\le q_c^2$ the entries obey
$E_{ck}^2\le\ell_c^{-1}\prod_{t=c}^{k-1}p_t^2/\ell_{t+1}$, which gives the forward recursion
$F=\sum_k\Phi_k$, $\Phi_k=(p_{k-1}^2\Phi_{k-1}+1)/\ell_k$, and the backward recursion
$R_c=(1+p_c^2R_{c+1})/\ell_c$ (exact-product form).
\end{lemma}

\begin{lemma}[RC chains]\label{lem:RC}
If the rows $R_{\rm lo},\dots,R_{\rm hi}$ are matched to the columns above them by $p$, then
$|E_{c,i}|=p_i^{-1}\prod_{t=i+1}^{c-1}q_t/p_t$ for $i<c$; if every $r_t\ge1$, then
$1/\lambda_{\min}\le\sum_i(\mathrm{hi}-i+1)/p_i^2$ and $R_c\le\sum_{i<c}1/p_i^2$.
\end{lemma}

\begin{lemma}[Tall chains]\label{lem:tall}
If the rows are $R_{\rm lo},\dots,R_{\rm hi}$, the columns $C_{\rm lo+1},\dots,C_{\rm hi}$,
and $r_t$ is strictly decreasing on the interior rows, then
$1/\lambda_{\min}(D^TD)\le\sum_{i<\mathrm{hi}}(\mathrm{hi}-i)/p_i^2+\sum_{i>\mathrm{lo}}(i-\mathrm{lo})/q_i^2$.
\end{lemma}

\begin{proof}
Delete the first row $j$ with $r_j\le1$ (or the top row). The lower piece is matched by $p$
with ratios $q/p\le1$, the upper by $q$ with ratios $p/q\le1$; add their Frobenius norms.
\end{proof}

In every bound the length factors ($\mathrm{hi}-c+1$, $\mathrm{hi}-i$, $i-\mathrm{lo}$) are
cancelled by a factor $n-L$ or $L-m$ of the links, which is what makes the bounds uniform.
The sums that remain are handled by the identity
$\sum_{1\le c\le k\le N}\frac1{k(N+1-c)}=H_N^{(2)}<\frac{\pi^2}6$, which follows from
grouping $\sum_{j+k\le N}1/(jk)$ by $j+k$.

\subsection{Worst cases and certificates}

Pair $1$ enters exactly. Pair $2$ enters only through monotone dependences ($P_2(L)$
increases in $n_2$ and decreases in $m_2$; $\Lambda^2$ is nonincreasing in $n_2^2$ and $m_2^2$
for $c\ge1$), so each regime is evaluated at the boundary of pair $2$: $n_2=n_1+\frac32$, and
$m_2$ at its extreme value; in R4 the pair-$1$ parameter is worst-cased instead. Every bound
is then a function of two parameters, $(n_1,m_1)$ or $(n_1,m_2)$, or of $n_1$ alone in R1c,
R4c, Tp and Tm. The hypotheses are $30$ exact polynomial certificates, each showing that a
polynomial in nonnegative shift variables has coefficients of one sign (or, once, that a
univariate polynomial has no root on $(1,\infty)$):
\begin{center}\small
\begin{tabular}{>{\raggedright\arraybackslash}p{1.3cm}>{\raggedright\arraybackslash}p{9.3cm}>{\raggedright\arraybackslash}p{2.6cm}}
\toprule
\estabhead{label} & \estabhead{claim} & \estabhead{used for}\\
\midrule
S1 & $\Lambda_c\le0$ on zigzag $0$ in R4 and R4c & $C_\Lambda =0$ there\\
S2 & $r$ strictly decreasing on interior rows of every tall chain & \cref{lem:tall}\\
S3 & $\bar r\le1$ for the CR chains of R1, R1c and R2$^\pm$ (zigzag $0$) & \cref{lem:CR}\\
S4 & $r\ge1$ for the RC chains of Tp and Tm & \cref{lem:RC}\\
S5 & $\partial\Lambda^2/\partial(n_2^2)\le0$, $\partial\Lambda^2/\partial(m_2^2)\le0$ for $c\ge1$ & worst cases\\
S6 & $\Lambda\le m_1^2[(4c+9)n_1^2+c(4c^2+11c+8)]/(4(2c+1)(2c+3))$ & tails\\
S7 & telescoping bounds $\bar r\le\frac{\Theta_c}{\Theta_c+1}\cdot(\dots)$, $\Theta_c=c-\mathrm{lo}+1$ & tails\\
S8 & two elementary rational inequalities & tails\\
\bottomrule
\end{tabular}
\end{center}
Negative controls (false variants of S1--S3) fail, as they must. The two columns $c=\frac12$
that occur (in R1c and in R2$^+$ with $m_1=\frac12$) are treated exactly:
$\Lambda_{1/2}=\frac5{16}(n_1^2-\frac94)-\frac13(n_2^2-4)\le0$ in R1c and
$\Lambda_{1/2}\le n_1^2/16$ in R2$^+$.

\subsection{The box}\label{subsec:box32}

For each parameter set the bound of \cref{cor:FW} is evaluated with $F_0,F_1,C_\Lambda$ the
smaller of the count and exact-product forms (\cref{lem:tall} for tall chains) and $\tau$ chosen
to balance the two terms. The parameter sets are all those with $n_1\le4000$ in the four
two-parameter regimes and $n_1\le10^5$ in the four one-parameter ones, $6.4\times10^7$ in
all. The minima of the resulting lower bound for $q=8\smin(B)$ are:
\begin{center}\small
\begin{tabular}{lrrl}
\toprule
\estabhead{regime} & \estabhead{parameter sets} & \estabhead{min bound for $q$} & \estabhead{at}\\
\midrule
R1 & $15\,980\,006$ & $10.873$ & $n_1=\frac72$, $m_1=\frac32$\\
R2$^+$ & $15\,988\,002$ & $7.125$ & $n_1=\frac92$, $m_1=\frac12$\\
R2$^-$ & $15\,984\,004$ & $4.723$ & $n_1=3$, $m_1=1$\\
R4 & $15\,999\,998$ & $6.507$ & $n_1=2$, $m_2=\frac32$\\
R1c & $99\,998$ & $8.219$ & $n_1=\frac52$\\
R4c & $99\,999$ & $4.621$ & $n_1=3$\\
Tp & $99\,999$ & $4.386$ & $n_1=4$\\
Tm & $99\,998$ & $3.777$ & $n_1=4$, i.e. $(\frac32;1,\frac32,\frac32,\frac32)$\\
\bottomrule
\end{tabular}
\end{center}
These values come from a binary64 evaluation, truncated to three decimals, and were
reproduced to six decimals by an independent exact-rational evaluation at the minimising
parameters. The certified version re-runs the whole box with every operation rounded outward
by one ulp in the safe direction and bounds the cancellation $1-\tau(\frac43+C_\Lambda)$ from
below directly. Every one of the $64\,352\,004$ parameter sets is certified positive, with no
NaN, and the certified lower bounds for the minima, truncated to six decimals, are
$10.873596$ (R1), $7.125868$ (R2$^+$), $4.723102$ (R2$^-$), $6.507913$ (R4), $8.219252$
(R1c), $4.621864$ (R4c), $4.386733$ (Tp) and $3.777422$ (Tm), each at the same parameters as
above. The smallest certified cancellation margin $1-\tau(\frac43+C_\Lambda)$ is
$3.02\times10^{-5}$ (truncated), at Tm. The binding value
$3.7774$ exceeds the constant $3.7$ of \cref{thm:z32}.

\subsection{Tails}

Beyond the box we use a sufficient condition that needs no choice of $\tau$.

\begin{lemma}\label{lem:sufficient}
Let $\lambda_*=(3.7/8)^2$ and $\mathcal C=\lambda_*[F_0+\frac73F_1+C_\Lambda F_1]$. If $\mathcal C\le1$, then
$\smin(B)^2\ge \lambda_*$.
\end{lemma}

\begin{proof}
Take $\tau=\theta/(\frac43+C_\Lambda)$ with $\theta=1-\lambda_*F_0$. The first term of \cref{cor:FW} is then
$\lambda_*$, and the second is at least $\lambda_*$ iff $\theta\ge \lambda_*(\frac43+C_\Lambda+\theta)F_1$, which
follows from $(1-\lambda_*F_0)(1-\lambda_*F_1)\ge1-\lambda_*F_0-\lambda_*F_1\ge \lambda_*(\frac43+C_\Lambda)F_1$, a consequence of
$\mathcal C\le1$.
\end{proof}

In each regime the three box quantities are bounded by closed forms in $n=n_1$ and
$\mathsf L=1+\ln n$, uniformly in the second parameter $a$. The factorisations behind them
are exact identities: for instance in R2$^-$ ($m_2=a-\frac12$, $\Theta_c=c-a+1$)
the squared links of zigzag $0$ satisfy
\[
  \frac{16}3\,\frac{c+1}c\,q_c^2=(n-c)(n+c+3)(\Theta_c+1)(c+1+a),
\]
and the same-column ratio is
\[
  \bar r_t=\frac{\Theta_t}{\Theta_t+2}\cdot\frac{t+1+a}{t+2+a}\cdot\frac{n+t+1}{n+t+4}\cdot\frac{(2t+5)(t+2)}{(2t+3)(t+1)},
\]
whose products telescope; $\Lambda$ is a quadratic in $n$ whose three coefficients are
bounded by certified polynomial inequalities. The resulting bounds are:
\begin{center}\small
\begin{tabular}{llll}
\toprule
\estabhead{regime} & \estabhead{$F_0$} & \estabhead{$F_1$} & \estabhead{$C_\Lambda$}\\
\midrule
R2$^-$ & $\frac{32}{3n}$ & $\frac{12}n+\frac{32\mathsf L}{3n(n+2)}$ & $8(2.5+\ln n)\mathsf L+\frac{400(2.5+\ln n)}{n-10}$\\[7pt]
R2$^+$ & $\frac{16}n$ & $\frac{32}{3n}+\frac{32\mathsf L}{3n(n+2)}$ & $16\mathsf L^2+\frac23\mathsf L$\\[7pt]
R1 & $\frac{4.22}n$ & $\frac{16}{3n}S(2a-3)$ & $a[2.67\mathsf L^2+12.6\mathsf L+9.3]$\\[7pt]
R4 & $\frac{16\pi^2}{18n}+\frac{64\mathsf L}{3n(n+3)}$ & $\frac{3.52}n+\frac{32\mathsf L}{3n(n+2)}$ & $0$\\[7pt]
R4c & $\frac{16\pi^2}{18n}+\frac{32\mathsf L}{3n(n+2)}$ & $\frac{16(\pi^2/6-1)}n+\frac{32\mathsf L}{3n(n+3)}$ & $0$\\[7pt]
R1c & $\frac{16(\mathsf L+\ln2)(1+2\mathsf L)}{3n^2}$ & $\frac{32\mathsf L^2}{n(n+1)}$ & $\frac{16}{225}(\mathsf L+\ln2)(1+2\mathsf L)$\\[7pt]
Tm & $\frac{32\mathsf L^2}{3n^2}$ & $\frac{16(\pi^2/6-1)}n+\frac{32\mathsf L}{3n(n+1)}$ & $6\mathsf L^2$\\[7pt]
Tp & $\frac{64\mathsf L^2}{3n^2}$ & $\frac{16(\pi^2/6-1)}n+\frac{32\mathsf L}{3n(n+2)}$ & $6\mathsf L^2$\\
\bottomrule
\end{tabular}
\end{center}
Here $S(x)=\sum_{i\ge1}1/(i(i+x))\le1/(1+x)+\ln(1+x)/x$, and in R1 the product
$C_\Lambda F_1$ is made uniform in $a$ by $aS(2a-3)\le3+\frac12\ln2n$. Writing $\mathcal C(n)$ as a sum
of monomials $K\mathsf L^p/\prod(n+\gamma_i)$, each is nonincreasing for $n\ge n_0$ once its
logarithmic derivative is negative at $n_0$, and $\mathcal C(n_0)$ is evaluated in $120$-bit
interval arithmetic, with every decimal constant entered as an exact rational and the
derivative test also done in intervals: $\mathcal C\le0.7644$ (R1), $0.7960$ (R2$^+$),
$0.5189$ (R2$^-$), $0.000915$ (R4) at $n_0=4000$, and $3\times10^{-6}$ (R1c),
$7.03\times10^{-5}$ (R4c), $0.0208$ (Tp, Tm) at $n_0=10^5$ (upper bounds, rounded up). All
are below $1$. (An earlier version of the script entered the decimal constants as binary64
numbers, so that its intervals enclosed the nearest binary64 points rather than the decimals,
and did the derivative test in ordinary precision. The margins are wide enough that its
verdicts were the same; the certificate of record is the repaired script,
\texttt{certificates/z32rigor/scripts/tail\_certify\_exact.py}.) The $50$ polynomial facts used are exact certificates,
and every closed form was also checked to dominate the corresponding box quantity at
$n_1\in\{60,60\frac12,200,200\frac12,500,999\frac12\}$ for all $a$.

\subsection{Small spins exactly}

Independently of box and tails, the two matrices of \eqref{eq:twochain} are tridiagonal
in the column index with rational diagonals and off-diagonals whose squares are rational,
so their positivity is decided by exact $LDL^T$ pivots at rational $\tau$. All $54\,019$ tuples
with $j_1=\frac32$ and spins at most $20$, counted as above, are certified this way at $q\ge7.18$; the smallest
certified value is $7.190$ at $(\frac32;\frac12,1,1,1)$, where the computed value is $7.476$, so the
two-chain inequality loses $4\%$ there.
